\documentclass[10pt]{article}
\usepackage[a4paper,margin=20mm]{geometry}
\usepackage{amsmath,amssymb,amsthm,mathtools,bm}
\usepackage{booktabs,array,microtype,xurl,hyperref}
\hypersetup{colorlinks=true,linkcolor=blue,citecolor=blue,urlcolor=blue}

\newtheorem{theorem}{Theorem}[section]
\newtheorem{proposition}[theorem]{Proposition}
\newtheorem{lemma}[theorem]{Lemma}
\newtheorem{corollary}[theorem]{Corollary}
\newtheorem{definition}[theorem]{Definition}
\newtheorem{remark}[theorem]{Remark}
\newcommand{\dd}{\,\mathrm d}
\newcommand{\RR}{\mathbb R}
\newcommand{\SSS}{\mathbb S}
\newcommand{\TT}{\mathbb T}
\newcommand{\cE}{\mathcal E}
\newcommand{\cF}{\mathcal F}
\newcommand{\cX}{\mathcal X}
\newcommand{\Tan}{\operatorname{Tan}}
\newcommand{\supp}{\operatorname{supp}}
\newcommand{\push}{_{\#}}

\title{AP-E16 Tangent Defects, Purified Caccioppoli Control,\\
and a Homogeneous Hopf-Base Young-Measure Obstruction}
\author{Route-F derivation ledger}
\date{24 July 2026}

\begin{document}
\maketitle

\begin{abstract}
This note executes the AP-E15 continuation without reopening the lattice
branch.  The complete-minor relaxation defect is split into five nonnegative
Hopf Radon measures.  At $\mu$-almost every point, one normalized tangent
sequence carries all five components with their Radon--Nikodym weights, and a
diagonal smooth recovery sequence can be chosen so that its normalized
compact-fibre-phase relaxation drop tends to zero.

A new scaling obstruction appears.  If the spatial radius is $r$, the target
amplitude is $s$, and first-, second-, and third-minor energies all contribute
to one defect mass, then their local scales are
\[
 s^2r,\qquad s^4/r,\qquad s^6/r^3.
\]
They admit one ordinary field blow-up only when $s\sim r$ and
$\mu(B_r)\sim r^3$.  Nonvolumetric tangents must be treated as order-specific
measure or generalized-Young-measure limits.

Inside one normal target ball, a degree-preserving geodesic cutoff gives a
conditional Caccioppoli inequality with exactly the AP-E13 weighted terms and
the normalized local comparison deficit.  It does not by itself yield excess
contraction.  Indeed an explicit smooth sphere-valued family, after exact
fibre-phase purification, generates the homogeneous pure-base Young measure
\[
 F_{31}=\cos\theta_1,\qquad F_{42}=\cos\theta_2.
\]
Its barycentric gradient and minors vanish, its weighted cutoff tends to
zero, but its diffuse graph-energy defect is
$\frac12+\frac R8>0$.  This disproves any unconditional inference
``phase purification plus two-limit weighted decay implies $\mu=0$''.  The
family is not asymptotically minimizing, so it does not refute the desired
theorem for the selected $B=1$ minimizer.  It proves that normalized full
local quasiminimality and a strong-quasiconvex boundary comparison are
indispensable.  Hessian, determinant, and portal promotion remain closed.
\end{abstract}

\tableofcontents

\section{Frozen energy and the Hopf defect vector}

The AP-E11 action remains
\begin{align}
 \cE(u)&=\int_\Omega\left[
 \frac12|Du|^2+\frac R2|M_2(Du)|^2+\frac K2|M_3(Du)|^2
 +m^2(1-u^0)\right]\dd x,\notag\\
 (R,K,m^2)&=(1,0.35,0.12).
 \label{eq:energy}
\end{align}
Let $u_j$ be a smooth boundary-preserving degree-$B$ sequence with
$\cE(u_j)\to I_B$, strongly convergent in $L^2$ to a graph representative
$u$, and weakly convergent at all three complete-minor levels.  AP-E15
defines
\begin{equation}
 e_j\dd x\stackrel{*}{\rightharpoonup}e_u\dd x+\mu,
 \qquad
 e_j=\frac12|Du_j|^2+\frac R2|M_2(Du_j)|^2
      +\frac K2|M_3(Du_j)|^2,\qquad \mu\geq0.
 \label{eq:mu}
\end{equation}

For $n_j=h(u_j)$ and $a_j=2u_j^*\theta$, set
\begin{align}
 \nu_j^a&=\frac18|a_j|^2\dd x,&
 \nu_j^n&=\frac18|Dn_j|^2\dd x,\notag\\
 \nu_j^{an}&=\frac R{32}|a_j\wedge Dn_j|^2\dd x,&
 \nu_j^c&=\frac R{32}|da_j|^2\dd x,\notag\\
 \nu_j^{cs}&=\frac K{128}|a_j\wedge da_j|^2\dd x.
 \label{eq:hopf-measures}
\end{align}
The exact AP-E14 identities imply
\begin{equation}
 e_j\dd x=\nu_j^a+\nu_j^n+\nu_j^{an}+\nu_j^c+\nu_j^{cs}.
 \label{eq:hopf-sum}
\end{equation}
Strong convergence of $u_j$ and weak graph convergence imply weak $L^2$
convergence of each Hopf factor to its counterpart for $u$.  Weighted lower
semicontinuity therefore gives, after one subsequence,
\begin{equation}
 \nu_j^\beta\stackrel{*}{\rightharpoonup}\nu_u^\beta+\mu_\beta,
 \qquad \mu_\beta\geq0,\qquad
 \mu=\sum_{\beta\in\{a,n,an,c,cs\}}\mu_\beta.
 \label{eq:component-defects}
\end{equation}

\begin{definition}[Hopf Radon--Nikodym weights]
Because every $\mu_\beta\leq\mu$, define
\begin{equation}
 \vartheta_\beta(x)=\frac{\dd\mu_\beta}{\dd\mu}(x),\qquad
 \vartheta_\beta\geq0,\qquad
 \sum_\beta\vartheta_\beta=1\quad\mu\text{-a.e.}
 \label{eq:weights}
\end{equation}
These weights distinguish fibre, base, shear, curvature, and
Chern--Simons-square defect.  Fibre-phase purification does not assert
$\vartheta_a=\vartheta_{an}=\vartheta_{cs}=0$; it removes only the part
exposed by an exact compact phase relaxation.
\end{definition}

\section[A common tangent and diagonal recovery]
{A common \(\mu\)-a.e. tangent and diagonal recovery}

Let $T_{x_0,r}(x)=(x-x_0)/r$.  At $\mu$-almost every $x_0$, standard
doubling-radius selection gives $r_\ell\downarrow0$ such that the normalized
measures
\begin{equation}
 \mu_{x_0,r_\ell}
 =\frac{(T_{x_0,r_\ell})_\#\mu}{\mu(B_{r_\ell}(x_0))}
 \stackrel{*}{\rightharpoonup}\tau\neq0
 \quad\text{locally in }\RR^3.
 \label{eq:tangent}
\end{equation}
The normalization can be chosen with
$\tau(B_1)=1$ after avoiding boundary-mass radii.

\begin{theorem}[Common Hopf tangent and phase-stationary diagonal]
For $\mu$-almost every $x_0$, the radii in \eqref{eq:tangent} may be chosen so
that
\begin{equation}
 \frac{(T_{x_0,r_\ell})_\#\mu_\beta}
      {\mu(B_{r_\ell}(x_0))}
 \stackrel{*}{\rightharpoonup}
 \vartheta_\beta(x_0)\tau
 \qquad\text{for every }\beta.
 \label{eq:component-tangent}
\end{equation}
There is also a diagonal $j_\ell\to\infty$ for which the scaled signed
recovery defects converge to the same measures and, for every fixed $L>0$,
\begin{equation}
 \frac{\Delta_{j_\ell}(B_{Lr_\ell}(x_0))}
      {\mu(B_{r_\ell}(x_0))}\longrightarrow0.
 \label{eq:normalized-phase-drop}
\end{equation}
Thus every tangent is stationary, in normalized energy, against compact
exact fibre shifts.
\end{theorem}

\begin{proof}
The Radon--Nikodym differentiation theorem for the finite family
$\mu_\beta=\vartheta_\beta\mu$ gives
\[
 \frac{\mu_\beta(B_r(x_0))}{\mu(B_r(x_0))}
 \longrightarrow\vartheta_\beta(x_0)
 \quad\text{for $\mu$-a.e. }x_0.
\]
Apply this simultaneously to a countable dense set of compactly supported
test functions after the doubling-radius selection.  This proves
\eqref{eq:component-tangent}.

Let $\sigma_j^\beta=\nu_j^\beta-\nu_u^\beta$.  For each already chosen
$r_\ell$, weak-star convergence permits $j_\ell$ to be so large that the
scaled errors against the first $\ell$ test functions are smaller than
$1/\ell$.  Since $\cE(u_j)-I_B\to0$, enlarge $j_\ell$ again so that
\begin{equation}
 \frac{\cE(u_{j_\ell})-I_B}{\mu(B_{r_\ell}(x_0))}
 <\frac1\ell.                                        \label{eq:diag-excess}
\end{equation}
AP-E15 proves
$0\leq\Delta_j(Q)\leq\cE(u_j)-I_B$ on every fixed admissible phase ball.
For fixed $L$, $Lr_\ell$ eventually lies below the phase-convexity radius.
Equations \eqref{eq:diag-excess} and the AP-E15 bound give
\eqref{eq:normalized-phase-drop}.
\end{proof}

\begin{remark}[Precise consequence of purification]
Equation \eqref{eq:normalized-phase-drop} is stronger than unnormalized
stationarity, but weaker than strong convergence of $a_j$.  A co-closed or
base-induced connection oscillation can be stationary along the exact fibre
orbit.  The tangent weights in \eqref{eq:weights} must therefore remain in
the ledger until a full local comparison eliminates them.
\end{remark}

\section{Why one field blow-up exists only in the volumetric branch}

Suppose a tangent is represented in a target chart by
\begin{equation}
 u(x_0+ry)=\exp_q\bigl(s\,v(y)\bigr),
 \label{eq:chart-blowup}
\end{equation}
where $r$ is the spatial scale and $s$ the target amplitude.  Up to uniformly
bounded chart factors, homogeneity of complete minors gives
\begin{align}
 \int_{B_r}|Du|^2\dd x
 &\sim s^2r\int_{B_1}|Dv|^2\dd y,\notag\\
 \int_{B_r}|M_2(Du)|^2\dd x
 &\sim \frac{s^4}{r}\int_{B_1}|M_2(Dv)|^2\dd y,\notag\\
 \int_{B_r}|M_3(Du)|^2\dd x
 &\sim \frac{s^6}{r^3}\int_{B_1}|M_3(Dv)|^2\dd y.
 \label{eq:three-scales}
\end{align}

\begin{proposition}[Simultaneous-normalization obstruction]
Assume $s=r^\alpha$ and $\mu(B_r)\asymp r^d$.  All three expressions in
\eqref{eq:three-scales} can be simultaneously of order $\mu(B_r)$ only if
\begin{equation}
 2\alpha+1=d,\qquad4\alpha-1=d,\qquad6\alpha-3=d,
 \label{eq:scaling-system}
\end{equation}
whose unique solution is
\begin{equation}
 \boxed{\alpha=1,\qquad d=3.}
 \label{eq:volumetric}
\end{equation}
\end{proposition}

\begin{proof}
The exponents in \eqref{eq:scaling-system} are the powers of $r$ in
\eqref{eq:three-scales}.  Subtracting the first equation from the second
gives $2\alpha=2$; hence $\alpha=1$, and then $d=3$.  The third equation is
consistent and gives no second branch.
\end{proof}

Thus a volumetric defect can have an ordinary tangent-gradient Young measure
with $s\sim r$.  At a nonvolumetric point, normalizing the first-order field
necessarily sends at least one higher-minor scale to zero or infinity.
The correct object is then the component tangent
\eqref{eq:component-tangent}, or a generalized Young measure with separate
concentration coordinates.  Treating every tangent as one $W^{1,2}$ field
would silently discard precisely the endpoint concentration under study.

\section{Conditional purified Hopf-base Caccioppoli inequality}

For a ball $B_t\Subset\Omega$, define the relative-homotopy local comparison
deficit
\begin{equation}
 \epsilon_j(B_t)=\cE(u_j;B_t)-
 \inf\left\{\cE(v;B_t):
 \begin{array}{l}
 v=u_j\text{ near }\partial B_t,\\
 v\simeq u_j\text{ relative }\partial B_t
 \end{array}\right\}.
 \label{eq:local-deficit}
\end{equation}
Replacing $u_j$ by an allowed $v$ in the full map preserves boundary data
and global degree.  Consequently
\begin{equation}
 0\leq\epsilon_j(B_t)\leq\cE(u_j)-I_B.
 \label{eq:deficit-global}
\end{equation}

Let the image of $u_j$ on $B_t$ lie in one fixed normal ball centred at
$q\in\SSS^3$, and put $\xi_j=\log_q u_j$.  For
$0<s<t$, choose $\eta=1$ on $B_s$, $\eta=0$ outside $B_t$, and
$|D\eta|\leq2/(t-s)$.  The comparison
\begin{equation}
 v_j=\exp_q\bigl((1-\eta)\xi_j\bigr)
 \label{eq:competitor}
\end{equation}
equals $q$ in $B_s$, equals $u_j$ near $\partial B_t$, and is homotopic to
$u_j$ relative to the boundary.

\begin{lemma}[Rank-one cutoff minor bounds]
With $A=D\xi_j$, $B=-\xi_j\otimes D\eta$, and $c=1-\eta$,
\begin{align}
 |D((1-\eta)\xi_j)|^2
 &\leq C(c^2|A|^2+|\xi_j|^2|D\eta|^2),\notag\\
 |M_2(cA+B)|^2
 &\leq C(c^4|M_2(A)|^2
          +c^2|\xi_j|^2|D\eta|^2|A|^2),\notag\\
 |M_3(cA+B)|^2
 &\leq C(c^6|M_3(A)|^2
          +c^4|\xi_j|^2|D\eta|^2|M_2(A)|^2).
 \label{eq:rank-one-bounds}
\end{align}
\end{lemma}

\begin{proof}
$B$ has rank one.  Hence
\[
 M_2(cA+B)=c^2M_2(A)+cL_2(A,B),\qquad
 M_3(cA+B)=c^3M_3(A)+c^2L_3(A,A,B),
\]
because every term containing two columns from $B$ vanishes.  Exterior
algebra gives
$|L_2|^2\leq C|A|^2|B|^2$ and
$|L_3|^2\leq C|M_2(A)|^2|B|^2$.  Together with
$|B|=|\xi_j||D\eta|$, this proves \eqref{eq:rank-one-bounds}.  The derivatives
of $\exp_q$ and its inverse only change $C$ on the fixed normal ball.
\end{proof}

Write
\begin{equation}
 \Phi_j(r)=\int_{B_r}\left[
 \frac12|Du_j|^2+\frac R2|M_2(Du_j)|^2
 +\frac K2|M_3(Du_j)|^2\right]\dd x.
 \label{eq:Phi}
\end{equation}

\begin{theorem}[Conditional Caccioppoli/hole-filling estimate]
There is a constant $C$ depending only on the normal ball and the frozen
coefficients such that
\begin{align}
 \Phi_j(s)
 \leq{}&C\bigl[\Phi_j(t)-\Phi_j(s)\bigr]\notag\\
 &+\frac{C}{(t-s)^2}\int_{B_t\setminus B_s}
 |\xi_j|^2\bigl(1+R|Du_j|^2+K|M_2(Du_j)|^2\bigr)\dd x\notag\\
 &+Cm^2\int_{B_t}|\xi_j|\dd x+\epsilon_j(B_t).
 \label{eq:caccioppoli}
\end{align}
Equivalently, with $\theta=C/(1+C)<1$,
\begin{align}
 \Phi_j(s)
 \leq{}&\theta\Phi_j(t)
 +\frac{C}{(t-s)^2}\int_{B_t\setminus B_s}
 |\xi_j|^2\bigl(1+R|Du_j|^2+K|M_2(Du_j)|^2\bigr)\dd x\notag\\
 &+Cm^2\int_{B_t}|\xi_j|\dd x+C\epsilon_j(B_t).
 \label{eq:hole-filling}
\end{align}
\end{theorem}

\begin{proof}
Use \eqref{eq:competitor} in \eqref{eq:local-deficit}.  The derivative energy
of $v_j$ vanishes on $B_s$.  Cancel the unchanged exterior, discard the
nonnegative derivative energy of $u_j$ on the annulus, and apply
\eqref{eq:rank-one-bounds}.  Because $V(u)=m^2(1-u^0)$ is $m^2$-Lipschitz on
$\SSS^3$,
\[
 \int_{B_t}|V(v_j)-V(u_j)|\dd x
 \leq Cm^2\int_{B_t}|\xi_j|\dd x.
\]
This proves \eqref{eq:caccioppoli}.  Moving
$C\Phi_j(s)$ to the left gives \eqref{eq:hole-filling}.
\end{proof}

After AP-E15 purification, the exact fibre-phase part of the normalized
comparison deficit vanishes.  The estimate above is therefore the legitimate
local inequality for the remaining base/mixed problem.  It still needs four
inputs before it becomes an excess theorem:
\begin{enumerate}
 \item normal-ball control or a trace-preserving endpoint replacement;
 \item sequence-uniform control of the weighted second- and third-order
       cutoff terms;
 \item normalized full local quasiminimality, not only phase stationarity;
 \item a strong-quasiconvex harmonic-approximation step which improves the
       scale exponent, rather than only $\theta<1$ hole filling.
\end{enumerate}

\section{A homogeneous phase-purified base Young measure}

The preceding list is not a technical luxury.  There is an exact obstruction
if the third input is omitted.  On $\TT^3=(\RR/2\pi\mathbb Z)^3$, for an
integer $N\geq2$, define
\begin{equation}
 u_N(x)=\left(
 \sqrt{1-\frac{\sin^2(Nx_1)+\sin^2(Nx_2)}{N^2}},\
 0,\
 \frac{\sin(Nx_1)}N,\
 \frac{\sin(Nx_2)}N
 \right).
 \label{eq:uN}
\end{equation}
It lies exactly in $\SSS^3$, converges uniformly to $e_0$, has rank at most
two, and is null-homotopic because its image lies in the open hemisphere
$u^0>0$.

\begin{proposition}[Positive diffuse pure-base defect]
The family \eqref{eq:uN} satisfies
\begin{align}
 Du_N&\rightharpoonup0,&
 M_2(Du_N)&\rightharpoonup0,&M_3(Du_N)&=0,\notag\\
 \frac1{|\TT^3|}\int_{\TT^3}|Du_N|^2\dd x&\longrightarrow1,&
 \frac1{|\TT^3|}\int_{\TT^3}|M_2(Du_N)|^2\dd x
 &\longrightarrow\frac14.
 \label{eq:uN-limits}
\end{align}
Its graph-energy measures converge to the nonzero diffuse defect
\begin{equation}
 e_N\dd x\stackrel{*}{\rightharpoonup}
 \left(\frac12+\frac R8\right)\dd x.
 \label{eq:diffuse-defect}
\end{equation}
Meanwhile its Hopf connection is
\begin{equation}
 a_N=\frac2N\left[
 -\sin(Nx_2)\cos(Nx_1)\dd x_1
 +\sin(Nx_1)\cos(Nx_2)\dd x_2\right],
 \label{eq:aN}
\end{equation}
so $\|a_N\|_2=O(N^{-1})$, while
\begin{equation}
 da_N=4\cos(Nx_1)\cos(Nx_2)\dd x_1\wedge\dd x_2.
 \label{eq:daN}
\end{equation}
The leading defect is therefore Hopf-base derivative/curvature, not an exact
fibre oscillation.
\end{proposition}

\begin{proof}
Write $\theta_i=Nx_i$.  The leading tangent derivatives are
\[
 \partial_1u_N=(0,0,\cos\theta_1,0)+O(N^{-1}),\qquad
 \partial_2u_N=(0,0,0,\cos\theta_2)+O(N^{-1}).
\]
All other columns vanish.  Periodic averaging gives
$\langle\cos\theta_i\rangle=0$,
$\langle\cos^2\theta_i\rangle=1/2$, and
$\langle\cos^2\theta_1\cos^2\theta_2\rangle=1/4$, proving
\eqref{eq:uN-limits}--\eqref{eq:diffuse-defect}.  Direct substitution in
$a=2u^*\theta$ gives \eqref{eq:aN}; differentiation gives \eqref{eq:daN}.
\end{proof}

Let $\chi_N$ minimize the periodic fibre-phase functional of $u_N$.  Since
$z_N=u_N^0>0$ is real, the potential contribution relative to $\chi=0$ is
$m^2z_N(1-\cos\chi)\geq0$.  Expanding every derivative term in
$a_N+2d\chi$ gives a coercive quadratic form in $d\chi$ plus a linear form
bounded by $C\|a_N\|_2\|d\chi\|_2$.  Comparison with $\chi=0$ therefore yields
\begin{equation}
 \|d\chi_N\|_2\leq C\|a_N\|_2=O(N^{-1}).
 \label{eq:phase-correction}
\end{equation}
Elliptic regularity makes $\chi_N$ smooth.  Thus
$\widetilde u_N=e^{i\chi_N}u_N$ is exactly phase-purified and has the same
leading defect.

The associated homogeneous tangent-gradient Young measure is the pushforward
of $(2\pi)^{-2}\dd\theta_1\dd\theta_2$ by
\begin{equation}
 F(\theta_1,\theta_2)=
 \begin{pmatrix}
 0&0&0\\
 0&0&0\\
 \cos\theta_1&0&0\\
 0&\cos\theta_2&0
 \end{pmatrix}.
 \label{eq:young-generator}
\end{equation}
It obeys
\begin{equation}
 \langle F\rangle=0,\qquad
 \langle M_2(F)\rangle=0,\qquad
 \left\langle\frac12|F|^2+\frac R2|M_2(F)|^2\right\rangle
 =\frac12+\frac R8.
 \label{eq:young-gap}
\end{equation}
It is therefore a nontrivial homogeneous complete-minor Young measure with
zero barycentric graph vector.

Most importantly,
\begin{equation}
 \frac1{|\TT^3|}\int_{\TT^3}
 |u_N-e_0|^2|M_2(Du_N)|^2\dd x
 =O(N^{-2})\longrightarrow0.
 \label{eq:weighted-zero}
\end{equation}
Hence even the AP-E15 two-limit weighted quantity can vanish while a diffuse
base defect remains positive.

\begin{theorem}[Exact scope of the no-go]
Phase purification, weak complete-minor graph convergence, and
sequence-uniform weighted cutoff decay do not imply $\mu=0$ or unconditional
normalized excess contraction.  The family $\widetilde u_N$ is a precise
counterexample.

It is not a counterexample to the desired theorem for the selected minimizer:
$u_N$ has degree zero and positive energy while the constant map has zero
energy.  Equivalently, its normalized full local comparison deficit does not
vanish.  The construction may be inserted into a vacuum region of a smooth
degree-one map without changing degree, but it still adds a removable
positive energy.
\end{theorem}

\section{Gate decision and next theorem}

AP-E16 changes the blocker in three ways:
\begin{enumerate}
 \item The $\mu$-a.e. common tangent, Hopf component weights, and normalized
       phase-stationary diagonal are \textbf{proved}.
 \item The geodesic-cutoff Caccioppoli inequality
       \eqref{eq:caccioppoli} is \textbf{proved conditionally} in one normal
       target ball, with every missing term displayed.
 \item Unconditional excess contraction from phase purification and weighted
       cutoff decay is \textbf{disproved} by the homogeneous base Young
       measure \eqref{eq:young-generator}.
\end{enumerate}

The next theorem must combine, rather than separate, recovery and
quasiminimality:
\begin{enumerate}
 \item On the volumetric branch \eqref{eq:volumetric}, prove a
       boundary-modification theorem which transfers normalized local
       quasiminimality to the tangent Young measure:
       \[
        \frac{\epsilon_{j_\ell}(B_{r_\ell})}
             {\mu(B_{r_\ell})}\longrightarrow0.
       \]
       Exact strong quasiconvexity would then force that measure to be Dirac
       and eliminate the diffuse branch.
 \item On nonvolumetric branches, work with the order-specific tangents
       \eqref{eq:component-tangent}; prove a concentration-capacity or
       topological lower-bound alternative rather than forcing one field
       normalization.
 \item Only after both branches are excluded may one assert full $\mu=0$,
       local recovery, regularity, and isolation.
\end{enumerate}

The bosonic Hessian remains unauthorized.  Dirac/Callias operators on a
prescribed smooth Hopf base may continue as parallel mathematics, but
determinant and degree-one portal promotion remain false.

\paragraph{Reproducibility.}
Run \texttt{python3}
\path{route_f/code/verify_ap_e16_tangent_caccioppoli_young_measure.py}
from the repository root.  The production card checks the scaling system,
rank-one minor algebra, homogeneous Young measure, exact sphere generator,
weighted-cutoff asymptotics, and all fail-closed gates.  It performs no
lattice relaxation or grid scan.

\end{document}
